\subsubsection{CRDT ბიბლიოთეკა (crdt-core)}

სისტემის მთავარი ნაწილია \texttt{crdt-core}, Rust-ზე დაწერილი დამოუკიდებელი
ბიბლიოთეკა, რომელიც YATA-ს სტილის ტექსტურ CRDT-ს მოიცავს. მისი პასუხისმგებლობაა დოკუმენტის
მდგომარეობა, ოპერაციების ინტეგრაცია და ბინარული (wire) ფორმატი.

\paragraph{ბლოკები და linked list.}
დოკუმენტი წარმოდგენილია \emph{ბლოკების} doubly linked list-ად. ბლოკი
არის ერთი ავტორის მიერ ზედიზედ აკრეფილი სიმბოლოების მონაკვეთი, რომელსაც
გააჩნია:
\begin{itemize}[itemsep=0pt]
  \item უნიკალური იდენტიფიკატორი \texttt{BlockId = (client\_id, clock)} -
        ავტორის ნომერი და მისი ლოკალური მთვლელი (ლამპორტის საათის
        ანალოგი).
        $n$ სიმბოლოიანი ბლოკი არის \texttt{(c, k)}-დან \texttt{(c, k+n-1)}-მდე;
  \item წარმოშობის ისტორია \texttt{origin\_left} და
        \texttt{origin\_right}, რომელი ორი ელემენტის შუაში დაიბადა ბლოკი
        ჩასმის მომენტში;
  \item შიგთავსი და წაშლის ბულეანი (\texttt{deleted}).
\end{itemize}

მთელი ტექსტის სიმბოლო-სიმბოლო შენახვის ნაცვლად ბლოკებით შენახვა
(მიდგომა, რომელიც Yjs-იდან გადმოვიღეთ \cite{yjs}) მეხსიერებასა და ქსელის
ტრაფიკს მკვეთრად ამცირებს: ჩვეულებრივი ბეჭდვისას მომდევნო სიმბოლო წინა
ბლოკს ეწებება/იმერჯება და ერთ ოპერაციად რჩება. ახალი ფაილის გახსნისას ტექსტი
64-სიმბოლოიან ბლოკებად იჭრება (\texttt{from\_text\_chunked}).

\paragraph{ბლოკის გაყოფა (splitting).}
რა ხდება, როდესაც ვინმე ბლოკის \emph{შუაში} ჩასვამს ტექსტს ან შუიდან
წაშლის? ბლოკი ორად იყოფა: მარცხენა ნახევარი ინარჩუნებს საწყის
იდენტიფიკატორს, მარჯვენა კი იღებს იმავე დიაპაზონის გაგრძელებას
(\texttt{(c, k+offset)}). რადგან იდენტიფიკატორები სინამდვილეში დიაპაზონებზეა
დამოკიდებული, ნებისმიერი
\texttt{BlockId} კვლავ ცალსახად პოულობს თავის ბლოკს, უბრალოდ ახლა შესაძლოა
ბლოკის შუაგულში აღმოჩნდეს. ეს ინვარიანტი (ე.წ. decomposition invariance)
კრიტიკულია: remote მონაწილის ოპერაცია შეიძლება ისეთ ბლოკზე განხორციელდეს,
რომელიც ჩვენთან უკვე სამ ნაწილად არის გაყოფილი.

\paragraph{ინტეგრაცია და კონფლიქტების გადაწყვეტა.}
remote ბლოკის მიღებისას ალგორითმმა უნდა იპოვოს მისი ზუსტი ადგილი
სიაში. თუ კონფლიქტი არ არის, ბლოკი უბრალოდ თავის origin-ებს შორის ჯდება.
თუ ორმა მონაწილემ ერთსა და იმავე ადგილას ერთდროულად ჩასვა ტექსტი
(ესეიგი თუ ორივე ბლოკს ერთი და იგივე origin-ები აქვს), YATA წესრიგს დეტერმინისტული
წესით ადგენს, კონფლიქტური ბლოკები ისე ლაგდება, რომ (1) ყველა რეპლიკაზე
ერთი და იგივე თანმიმდევრობა მიიღება და (2) ერთი ავტორის მიმდევრობითი
ჩანაწერები ერთად რჩება. Rust-ის type სისტემა აქ
ძალიან დაგვეხმარა: ინტეგრაციის ლოგიკა მკაცრად გამიჯნულ მოდულებშია
(\texttt{document/integrate.rs}) და ყოველი შემთხვევა ცალკე
ტესტითაა დაფარული.

\paragraph{წაშლა, DeleteSet და state ვექტორი.}
CRDT-ში წაშლა განსხვავებულად მუშაობს: ბლოკის ფიზიკურად ამოღება არ
შეიძლება, რადგან სხვა მონაწილის ჯერ მოუსვლელი ოპერაცია შესაძლოა სწორედ
მას ეყრდნობოდეს origin-ად. ამიტომ წაშლა მხოლოდ \emph{მონიშვნაა}: ბლოკი
რჩება სიაში როგორც tombstone, ტექსტში კი აღარ ჩანს.
წაშლები კომპაქტურად ინახება \texttt{DeleteSet} სტრუქტურაში, თითოეული
კლიენტისთვის წაშლილი \texttt{clock}-დიაპაზონების სიად.

\emph{state ვექტორი} ინახავს თითო კლიენტისთვის მისგან ნანახი ოპერაციების
მაქსიმალურ \texttt{clock}-ს. ორი რეპლიკის ვექტორების შედარებით ზუსტად ჩანს, რომელ
მხარეს აკლია რომელი ოპერაციები. ჩვენ ვექტორებს ორ რამეში ვიყენებთ:
გვიან შემოსული სტუმრის სინქრონიზაციასა და garbage collector-ში (იხ.
ოპტიმიზაციების თავი), სადაც ის განსაზღვრავს, რომელი წაშლა აქვს ყველას
გარანტირებულად ნანახი.

\paragraph{სნეფშოტები და wire პროტოკოლი.}
დოკუმენტის სრული მდგომარეობა, ბლოკები, DeleteSet, state ვექტორი -
სერიალიზდება \texttt{Snapshot} სტრუქტურად (ბინარულად, \texttt{bitcode}
ბიბლიოთეკით). ქსელში ყოველ ფრეიმს ერთბაიტიანი პრეფიქსი აქვს, რომელიც
ტიპს განსაზღვრავს: ოპერაცია (\texttt{0x00}), სნეფშოტი (\texttt{0x01}),
საკონტროლო ფრეიმი (\texttt{0x02}), garbage collector commit
(\texttt{0x04}), წევრობის ცვლილება (\texttt{0x05}), state ვექტორის
რეპორტი (\texttt{0x06}), უფლების ცვლილება (\texttt{0x07}) და მონაწილის
ინფორმაცია (\texttt{0x08}). რადგან ამ ფორმატს ორი ენა კითხულობს (Rust
კლიენტებში, Go გეითვეიში), ორივე მხარეს გვაქვს ე.წ. protocol drift
ტესტები.

\paragraph{პოზიციური B+ ინდექსი.}
CRDT ოპერაციები იდენტიფიკატორებზე მუშაობს, რედაქტორი კი, სიმბოლოს
პოზიციებზე („ჩასვი 42-ე პოზიციაზე“). ამ ორ სამყაროს შორის გადაყვანა
linked list-ის გავლით $O(n)$ ოპერაციაა და ყოველ კლავიშზე მთელი
დოკუმენტის ჩამოვლა დიდ ფაილებზე შესამჩნევად ანელებდა სისტემას. ამის
გადასაჭრელად დოკუმენტს დავურთეთ \emph{აუგმენტირებული B+ ხე}: ბლოკები ხის
ფოთლებში დოკუმენტის თანმიმდევრობისავით აწყვია, ხოლო ყოველი შიდა node
ინახავს შვილი ნოუდის (ქვე ხის) ხილული სიმბოლოების ჯამურ სიგრძეს (\texttt{visible\_len};
წაშლილი ბლოკის ხილული სიგრძე ნულია). პოზიციით ძებნისას ხეზე დაშვებისას
თითო დონეზე ვაკლებთ მარცხენა შვილების ჯამებს, $O(\log n)$; პირიქით,
ბლოკიდან პოზიციის გამოთვლისას ფოთლიდან ზემოთ ავდივართ და მარცხენა ძმების
ჯამებს ვაგროვებთ.

მნიშვნელოვანი დეტალი: ხე \emph{მეორეული, წარმოებული} სტრუქტურაა -
ჭეშმარიტების წყარო ყოველთვის doubly linked list რჩება. იმისთვის, რომ ისინი
არასდროს აცდნენ ერთმანეთს, debug რეჟიმში ყოველი მუტაციის შემდეგ ეშვება
ფუნქცია, რომელიც მთელ სიას ხელით ჩამოუვლის და ამოწმებს, რომ ხის პასუხები
ზუსტად ემთხვევა. 
ხის child node size განზრახ განსხვავებულია: debug-ში მცირეა (4),
რომ ტესტებში გაყოფის ლოგიკა ხშირად გავარჯიშდეს, release-ში კი დიდი რიცხვი ავიღეთ
(64/32), ნაკლები სიღრმე და უკეთესი კეშ-ლოკალურობა.

\begin{diagram}[htbp]
  \centering
  \begin{tikzpicture}[
      font=\small\sffamily,
      blk/.style={draw=pcgreen, thick, minimum height=1.0cm, minimum width=2.1cm,
                  align=center, fill=pcgreenL, rounded corners=2pt},
      tomb/.style={blk, draw=pcred, fill=pcredL},
      tnode/.style={draw=pcblue, thick, rounded corners=3pt, align=center,
                    fill=pcblueL, minimum height=0.9cm, inner sep=5pt},
      arr/.style={-{Stealth}, very thick, draw=black!70}
    ]
    % root
    \node[tnode] (root) at (4.5,3.4) {\textbf{root}\\{[}left: 9 vis{]} {[}right: 5 vis{]}};
    % leaves
    \node[tnode] (leafA) at (1.3,1.8) {leaf 1};
    \node[tnode] (leafB) at (7.5,1.8) {leaf 2};
    \draw[arr, draw=pcblue] (root) -- (leafA);
    \draw[arr, draw=pcblue] (root) -- (leafB);

    % linked list
    \node[blk]  (b1) at (0,0)   {(1,0)\\``Hel''};
    \node[blk]  (b2) at (2.4,0) {(1,3)\\``lo ''};
    \node[tomb] (b3) at (4.8,0) {(2,0)\\\emph{del.} ``XY''};
    \node[blk]  (b4) at (7.2,0) {(2,2)\\``wor''};
    \node[blk]  (b5) at (9.6,0) {(1,6)\\``ld''};
    \draw[arr] (b1) -- (b2);
    \draw[arr] (b2) -- (b3);
    \draw[arr] (b3) -- (b4);
    \draw[arr] (b4) -- (b5);

    \draw[arr, dashed, draw=pcblue!70] (leafA) -- (b1);
    \draw[arr, dashed, draw=pcblue!70] (leafA) -- (b2);
    \draw[arr, dashed, draw=pcblue!70] (leafB) -- (b3);
    \draw[arr, dashed, draw=pcblue!70] (leafB) -- (b4);
    \draw[arr, dashed, draw=pcblue!70] (leafB) -- (b5);
  \end{tikzpicture}
  \caption{CRDT დოკუმენტის სტრუქტურა: ბლოკების linked list
  (ქვემოთ, tombstone წითლადაა) და მასზე აგებული პოზიციური B+ ინდექსი
  (ზემოთ, კვანძებში, ხილული სიმბოლოების ჯამები)}
  \label{diag:crdt}
\end{diagram}

ჩასმისა და წაშლის დინამიკას ორი კადრ-კადრ დაშლილი დიაგრამა აჩვენებს
(დიაგრამები~\ref{diag:insert} და~\ref{diag:delete}): პირველი, თუ როგორ
ჯდება ახალი ბლოკი ერთდროულად სიაშიც და ხის ფოთოლშიც და როგორ „ბუბლდება“
ხილული სიგრძის დელტა ფესვამდე; მეორე, თუ როგორ იქცევა ბლოკი
tombstone-ად და როგორ სწორდება ჯამები მთელ შტოზე.

\begin{diagram}[htbp]
  \centering
  \begin{tikzpicture}[
      font=\footnotesize\sffamily,
      blk/.style={draw=pcgreen, thick, minimum height=0.85cm, minimum width=1.7cm,
                  align=center, fill=pcgreenL, rounded corners=2pt},
      newblk/.style={blk, draw=pcorange, fill=pcorangeL, very thick},
      tnode/.style={draw=pcblue, thick, rounded corners=3pt, align=center,
                    fill=pcblueL, minimum height=0.75cm, inner sep=4pt},
      hot/.style={draw=pcorange, very thick},
      arr/.style={-{Stealth}, thick, draw=black!70},
      steplbl/.style={anchor=west, font=\small\bfseries, text=pcblue}
    ]
    % ------------- Frame 1 -------------
    \begin{scope}[yshift=0cm]
      \node[steplbl] at (-0.9,3.0) {ნაბიჯი 1, საწყისი მდგომარეობა};
      \node[tnode] (r1) at (2.6,2.2) {root: [5 vis]};
      \node[tnode] (l1) at (2.6,1.1) {leaf: [(1,0):3] [(1,3):2]};
      \draw[arr, draw=pcblue] (r1) -- (l1);
      \node[blk] (a1) at (1.2,-0.1) {(1,0)\\``Hel''};
      \node[blk] (b1) at (4.0,-0.1) {(1,3)\\``lo''};
      \draw[arr] (a1) -- (b1);
    \end{scope}
    % ------------- Frame 2 -------------
    \begin{scope}[yshift=-4.3cm]
      \node[steplbl] at (-0.9,3.0) {ნაბიჯი 2, ახალი ბლოკი ჯდება სიაშიც და ფოთოლშიც};
      \node[tnode] (r2) at (2.6,2.2) {root: [5 vis]};
      \node[tnode, hot] (l2) at (2.6,1.1) {leaf: [(1,0):3] \textbf{[(2,0):2]} [(1,3):2]};
      \draw[arr, draw=pcblue] (r2) -- (l2);
      \node[blk] (a2) at (0.4,-0.1) {(1,0)\\``Hel''};
      \node[newblk] (n2) at (2.6,-0.1) {(2,0)\\``XY''};
      \node[blk] (b2) at (4.8,-0.1) {(1,3)\\``lo''};
      \draw[arr] (a2) -- (n2);
      \draw[arr] (n2) -- (b2);
    \end{scope}
    % ------------- Frame 3 -------------
    \begin{scope}[yshift=-8.6cm]
      \node[steplbl] at (-0.9,3.0) {ნაბიჯი 3, დელტა (+2) ბუბლდება ფესვამდე};
      \node[tnode, hot] (r3) at (2.6,2.2) {root: [\textbf{7} vis]};
      \node[tnode] (l3) at (2.6,1.1) {leaf: [(1,0):3] [(2,0):2] [(1,3):2]};
      \draw[arr, draw=pcblue] (r3) -- (l3);
      \draw[-{Stealth}, very thick, draw=pcorange]
        ($(l3.east)+(0.25,0)$) to[bend right=40]
        node[right, font=\small\sffamily\bfseries, text=pcorange]{+2} ($(r3.east)+(0.25,0)$);
      \node[blk] (a3) at (0.4,-0.1) {(1,0)\\``Hel''};
      \node[blk, draw=pcorange] (n3) at (2.6,-0.1) {(2,0)\\``XY''};
      \node[blk] (b3) at (4.8,-0.1) {(1,3)\\``lo''};
      \draw[arr] (a3) -- (n3);
      \draw[arr] (n3) -- (b3);
    \end{scope}
  \end{tikzpicture}
  \caption{ბლოკის ჩასმა ნაბიჯ-ნაბიჯ: ბლოკი ერთდროულად ემატება
  list-ს და ხის ფოთოლს, შემდეგ კი ხილული სიგრძის დელტა
  ფესვამდე ვრცელდება}
  \label{diag:insert}
\end{diagram}

\begin{diagram}[htbp]
  \centering
  \begin{tikzpicture}[
      font=\footnotesize\sffamily,
      blk/.style={draw=pcgreen, thick, minimum height=0.85cm, minimum width=1.7cm,
                  align=center, fill=pcgreenL, rounded corners=2pt},
      tomb/.style={blk, draw=pcred, fill=pcredL},
      tnode/.style={draw=pcblue, thick, rounded corners=3pt, align=center,
                    fill=pcblueL, minimum height=0.75cm, inner sep=4pt},
      hot/.style={draw=pcred, very thick},
      arr/.style={-{Stealth}, thick, draw=black!70},
      steplbl/.style={anchor=west, font=\small\bfseries, text=pcblue}
    ]
    % ------------- Frame 1 -------------
    \begin{scope}[yshift=0cm]
      \node[steplbl] at (-0.9,3.0) {ნაბიჯი 1, ბლოკი ინიშნება წაშლილად (tombstone)};
      \node[tnode] (r1) at (2.6,2.2) {root: [7 vis]};
      \node[tnode, hot] (l1) at (2.6,1.1) {leaf: [(1,0):3] \textbf{[(2,0):0]} [(1,3):2]};
      \draw[arr, draw=pcblue] (r1) -- (l1);
      \node[blk] (a1) at (0.4,-0.1) {(1,0)\\``Hel''};
      \node[tomb] (n1) at (2.6,-0.1) {(2,0)\\\emph{del.} ``XY''};
      \node[blk] (b1) at (4.8,-0.1) {(1,3)\\``lo''};
      \draw[arr] (a1) -- (n1);
      \draw[arr] (n1) -- (b1);
    \end{scope}
    % ------------- Frame 2 -------------
    \begin{scope}[yshift=-4.3cm]
      \node[steplbl] at (-0.9,3.0) {ნაბიჯი 2, დელტა ($-2$) ბუბლდება ფესვამდე};
      \node[tnode, hot] (r2) at (2.6,2.2) {root: [\textbf{5} vis]};
      \node[tnode] (l2) at (2.6,1.1) {leaf: [(1,0):3] [(2,0):0] [(1,3):2]};
      \draw[arr, draw=pcblue] (r2) -- (l2);
      \draw[-{Stealth}, very thick, draw=pcred]
        ($(l2.east)+(0.25,0)$) to[bend right=40]
        node[right, font=\small\sffamily\bfseries, text=pcred]{$-2$} ($(r2.east)+(0.25,0)$);
      \node[blk] (a2) at (0.4,-0.1) {(1,0)\\``Hel''};
      \node[tomb] (n2) at (2.6,-0.1) {(2,0)\\\emph{del.} ``XY''};
      \node[blk] (b2) at (4.8,-0.1) {(1,3)\\``lo''};
      \draw[arr] (a2) -- (n2);
      \draw[arr] (n2) -- (b2);
    \end{scope}
  \end{tikzpicture}
  \caption{ბლოკის წაშლა ნაბიჯ-ნაბიჯ: ბლოკი სიაში რჩება tombstone-ად,
  ხილული სიგრძე ნულდება და უარყოფითი დელტა ხის ჯამებს ასწორებს}
  \label{diag:delete}
\end{diagram}
